\section{Discussion and limitations}\label{sec:discussion}
The experiments establish timing-only vulnerability in the implemented systems and show why the threat model and computational comparison are part of the result. The hybrid study identifies constrained retimings that change final predictions, while the classical grid study exhibits a large contrast under common local conditions. Their different samples, victims, representations and budget units prevent a single robustness ranking across the two studies.

The controlled architecture gap is the strongest observed contrast, but its mechanism remains unresolved. The integer custom SNN and spiking ResNet18 have nearly equal clean accuracy and radically different ASR, motivating direct investigation of temporal dynamics, decision geometry and architecture-dependent search difficulty. Shared nominal settings do not guarantee equally effective optimization. The low-ASR controls require stronger or adaptive attacks and independently trained seeds before a broader robustness conclusion is warranted.

The study does not isolate a quantum-specific failure mechanism. QSNN-v3 has a classical frontend and an analytically simulated circuit, and final TEMP-DRIFT-v2 optimizes class margin rather than quantum-state drift. Historical Iris diagnostics, retained in the supplement, show that greater mean state drift need not coincide with more classification failures. A mechanism study should relate final-model state distance to per-example outcomes while controlling clean margin and realized input distortion.

\paragraph{Encoding-aware interpretation.}
The quantum component is only one stage of the attacked composition. A feasible timestamp perturbation can change temporal aggregation, spiking summaries, encoded angles and the final probability vector. Without interventions that hold the frontend and training protocol fixed while replacing the readout, the observed failures cannot be attributed uniquely to quantum processing. This distinction complements comparisons of quantum and classical decision geometry \cite{wendlinger2024comparative} and avoids transferring guarantees between incompatible encoding or attacker-access settings \cite{gong2024randomized,dowling2026guarantees}.

\paragraph{Optimization diagnostics.}
The hard-forward, surrogate-backward timestamp pipeline makes attack diagnostics particularly relevant. Gradient checks, alternative surrogates, multiple initializations and gradient-free searches would help distinguish optimizer failure from difficult decision boundaries \cite{athalye2018obfuscated}. The principle of diverse attack evaluation \cite{croce2020reliable} is applicable, but an image-space ensemble cannot be imported unchanged: every candidate must obey the native event constraints or the packet-grid capacity rules. The present candidate and time controls are useful computational comparisons, not evidence that the search has exhausted either feasible set.

The remaining boundaries are concrete. Five hybrid seeds share one frontend and split; timestamp attacks reuse a selected validation population and one victim seed; grid attacks use deterministic joint-correct subsets. Official-test QSNN accuracy and robustness are absent. Tomography monitoring has no final N-MNIST detection or defense evaluation. This manuscript verifies report arithmetic rather than rerunning models or attacks, and original audit verdicts are attributed to the preserved reports because most underlying artifacts are missing.

\paragraph{Event preservation and stealth are distinct.}
An unchanged event count does not imply unchanged inter-spike intervals, frame histograms or detector behavior. The final-model experiments contain no statistical-indistinguishability test or adaptive attack against a deployed monitor. They are test-time evasion experiments; they do not modify a training set.

\paragraph{Next empirical steps.}
Future evaluation should select the final model using validation, then measure QSNN clean accuracy and attacks on the official test partition with independently trained frontend and head seeds, sample identifiers and each model's clean errors. Attack diagnostics should sweep restarts, iteration counts and step sizes, cross-entropy and margin losses, alternative surrogates and derivative-free budget curves. The best valid success found up to each budget should be recorded and reconfirmed with the hard forward pass.

A future component ablation should hold the frontend, split and training budget fixed while comparing a capacity-matched classical head and variational heads with and without entangling gates, reporting parameter counts and training costs. This ablation has not been performed. A state/margin study should retain each example's final state, probability vector and clean and attacked margins, comparing failures and survivors while controlling realized timestamp distortion and clean margin. Aggregate mean fidelity alone would not establish a mechanism.

The historical corrected integer campaign reports only 8 of 33 conditions complete. A proposed tomography monitor must first specify eight-qubit observables, a shot budget, clean-only calibration and a rejection or recovery policy, then measure ROC/PR curves, false-positive rates and adaptive-attack outcomes. These final-model defense results are unavailable. Local single-qubit marginals do not determine the full joint state. The supplementary mathematical analysis and earlier Iris diagnostics inform that proposal without supplying a final-model defense result.

\paragraph{Evidence availability.}
The comparison report contains results although its accompanying README states that no experiment results are present. We use the populated report and disclose this inconsistency. Most original checkpoints, raw per-example arrays, attack manifests and logs are absent from the accessible snapshot. We therefore distinguish checked report arithmetic from independently reproduced model runs and attribute the historical audit verdicts to their reports.

\section{Conclusion}\label{sec:conclusion}
QSNN-v3 provides an implemented eight-qubit hybrid victim with 97.272\% $\pm$ 0.212\% validation accuracy, and constrained timestamp retiming changes its decisions on 3--5 of the 100 examined joint-correct validation examples. Paired comparisons establish no significant advantage of TEMP-DRIFT-v2 over timestamp PGD; its largest numerical gain occurs under elapsed-time matching. In the separate common-recipe packet-grid study, the custom SNN exhibits 98.0--100\% ASR while three classical reference architectures exhibit 0--1.9\% under the two completed budgets. The resulting contribution is an explicit account of measurable retiming vulnerability within specified computational and representational settings. Official-test QSNN robustness, a quantum-specific failure mechanism and an effective final monitoring defense remain open empirical questions. The resulting evidence supports a protocol-specific account of end-to-end temporal vulnerability. It also identifies what stronger mechanistic claims would require: independent test evaluation, architecture-aware attack diagnostics, and per-example links between input distortion, quantum-state displacement and classification margin.

\section*{Data and code availability}
Implementation and preserved reports are available in the \texttt{ABC\_def1} repository \cite{repository}. The supplied supplement and data tables identify original report paths and the pinned snapshot. Tables transcribe available reports; they are not newly generated model outputs. Full experimental reproduction requires the omitted artifacts and original datasets. Author details are withheld for review, consistent with the supplied manuscript.
